\section{Recovering curve subfields}
\label{sec:curve-fields}

The support bound of Proposition~\ref{prop:uniform-support} allows us to
pass from individual multiplicative classes to subfields.  We prove that
$\Theta$ carries the mod-$\ell$ multiplicative group of every relatively
algebraically closed one-variable subfield of $K$ onto that of such a
subfield of $L$.  The geometric step is an incidence construction: infinitely
many divisors of bounded degree form a family, and a curve in its parameter
space produces the one-variable field.  This follows the bounded-support
and incidence method of \cite[Sections~4--6]{InternalBP}.  We give the
construction and the descent argument in the finite-coefficient setting.

A \emph{curve subfield} of a function field $F/\kappa$ is a subfield
$E\subset F$ containing $\kappa$, relatively algebraically closed in $F$,
and of transcendence degree one over $\kappa$.  It is finitely generated:
for any $t\in E\setminus\kappa$, it is the relative algebraic closure of
$\kappa(t)$ in $F$, which is finite over $\kappa(t)$.  The natural map
$V_E\to V_F$ is injective, since an $\ell$-th root in $F$ of an element
of $E$ is algebraic over $E$.  We henceforth regard $V_E$ as a subspace
of $V_F$.

\subsection{Two descent facts}

We first record the field-theoretic facts needed to descend the field
produced by incidence.  Recall that a finitely generated extension $F/B$
is \emph{regular} if it is separable and $B$ is relatively algebraically
closed in $F$; equivalently, $F$ and an algebraic closure $\overline B$ are
linearly disjoint over $B$.

\begin{lemma}[Descent from algebraically closed constants]
\label{lem:constant-descent}
Let $F/B$ be a finitely generated regular extension, where
$\charac(B)\ne\ell$ and $\mu_\ell\subset B$.  If $g\in F^\times$
has an $\ell$-th root in $F\overline B$, then its class in $V_F$ lies
in the image of $V_B$.
\end{lemma}

\begin{proof}
Choose $r\in F\overline B$ with $r^\ell=g$.  The element $r$ already
belongs to $FB'$ for a finite Galois extension $B'/B$.  Indeed it is
separable over $F$, so purely inseparable constant extensions are
unnecessary, and we may take a finite Galois closure of the remaining
constant extension.  By regularity,
$\Gal(FB'/F)=\Gal(B'/B)$.  For $\gamma$ in this group,
\[
                         c_\gamma=\frac{\gamma(r)}r\in\mu_\ell.
\]
These multipliers form a multiplicative cocycle with values in $B'^\times$.
Hilbert's Theorem~90 gives $b'\in B'^\times$ with
$\gamma(b')/b'=c_\gamma$ for every $\gamma$.  Thus $r/b'\in F$ and
$b'^\ell\in B$.  The equality
$g=(r/b')^\ell b'^\ell$ proves the assertion.
\end{proof}

\begin{lemma}[Intersections of curve subfields]
\label{lem:curve-intersection}
Let $F/\kappa$ be a finitely generated extension of an algebraically
closed field of characteristic different from $\ell$.  If $P,Q\subset F$
are distinct curve subfields, then
\[
                          \dim_\Lambda(V_P\cap V_Q)<\infty.
\]
Each $V_P$ is infinite dimensional.  In particular, an inclusion
$V_P\subset V_Q$ forces $P=Q$.
\end{lemma}

\begin{proof}
The compositum $PQ$ has transcendence degree two over $\kappa$.
Otherwise every element of $Q$ would be algebraic over $P$, so relative
algebraic closedness of $P$ in $F$ would give $Q\subset P$; reversing
their roles would give equality.
Let $C_P,C_Q$ be their smooth projective curves.  The natural map to
$C_P\times C_Q$ is dominant, since its image has dimension two;
consequently $PQ$ is the function field of this product.

Inside $V_{PQ}$, any class coming from both fields has zero order along
every divisor $\{x\}\times C_Q$.  Its representative from $P$ therefore
lies in the common kernel of all point orders of $C_P$.  This kernel has
dimension $2g(C_P)$ by Lemma~\ref{lem:curve-orders}, so the intersection
inside $V_{PQ}$ is finite dimensional.  The maps from $V_P$ and $V_Q$
to $V_{PQ}$ are injective: each factor field is relatively algebraically
closed in the function field of the product.

The kernel of $V_{PQ}\to V_F$ is finite dimensional by
Lemma~\ref{lem:kummer-kernel}.  To see that passing to $F$ preserves the
finite-intersection conclusion, consider pairs $(a,b)\in V_P\oplus V_Q$
whose images in $V_F$ agree.  The difference $a-b$ lies in this finite
kernel; the kernel of the difference map on such pairs is the intersection
already computed in $V_{PQ}$.  The space of these pairs, and hence
$V_P\cap V_Q$ in $V_F$, is finite dimensional.

Finally, choose $t\in P\setminus\kappa$.  The classes $[t-a]$,
$a\in\kappa$, are linearly independent in $V_{\kappa(t)}$, as their
orders at the distinct finite points show.  Lemma~\ref{lem:kummer-kernel}
gives only a finite-dimensional kernel on passing to $V_P$, so $V_P$
is infinite dimensional.
\end{proof}

\subsection{A family of divisors and a field of constants}

The next lemma isolates the geometric construction.  Its hypothesis says
that every class under consideration becomes an $\ell$-th power on almost
every divisor in one fixed infinite family.  The conclusion realizes all
these classes as classes from a single curve field after a finite extension
of the ambient field.

\begin{lemma}[Incidence descent]
\label{lem:curve-incidence}
Let $X\subset\Pj_l^n$ be a normal integral projective variety of dimension
$d\ge2$ over an algebraically closed field $l$ of characteristic different
from $\ell$, and put $L=l(X)$.  Let $\mathcal S$ be an infinite set of
prime divisors on $X$ whose degrees are bounded.  Suppose that
$H\subset V_L$ is a subspace with the following property: for every
$g\in L^\times$ with $[g]\in H$,
\begin{equation}
\label{eq:curve-residue-trivial}
 \ord_D(g)=0\quad\hbox{and}\quad
 g|_D\in l(D)^{\times\ell}
 \qquad\hbox{for all but finitely many }D\in\mathcal S.
\end{equation}
Then there are a finite extension $M/L$ and a one-variable field
$P_0/l$ contained in $M$ such that $M/P_0$ is regular,
\begin{equation}
\label{eq:curve-compositum}
                               M=LP_0,
\end{equation}
and the image of $H$ in $V_M$ is contained in the image of $V_{P_0}$.
In particular, every $g$ as above has an $\ell$-th root in
$M\overline{P_0}$.
\end{lemma}

\begin{proof}
\emph{A bounded parameter space.}
We recall why a degree bound gives a parameter space of finite type.
Integral subvarieties of fixed dimension
and bounded degree in a fixed projective space have only finitely many
Hilbert polynomials; see also \cite[Proposition~5.3]{Tanaka2024}.
One way to establish this boundedness is to cut them
out set-theoretically by forms of bounded degree.  If an integral
subvariety $D\subset\Pj^n$ has dimension $r$ and degree at most $b$, and
$x\notin D$, choose a linear projection to $\Pj^{r+1}$ whose center
avoids the join of $x$ and $D$.  When $D$ is already a hypersurface no
projection is necessary.  The projection is defined on $D$ and is finite
there: a positive-dimensional fiber would contradict ampleness of the
pulled-back hyperplane bundle.  Its image is a hypersurface of degree at
most $b$, and its equation pulls back to a form vanishing on $D$ but not
at $x$.  Multiplication by a form nonzero at $x$ makes the degree exactly
$b$ if necessary.  It follows that the degree-$b$ forms vanishing on $D$
cut it out as a reduced set.

A fixed number of such forms, namely the dimension of the space of
degree-$b$ forms, therefore places all these reduced subvarieties among
the geometric reductions of fibers of one projective family of finite
type.  Such geometric reductions have finitely many Hilbert polynomials.
For completeness, over an integral base descend the reduction of the
geometric generic fiber to a finite extension of the base function field,
and normalize the base in that extension.  On a dense open of this finite
base change, spread the reduction as a closed subscheme, with nilpotent
defining ideal, flat and with geometrically reduced fibers.  These
properties follow after shrinking from generic flatness and spreading
geometric reducedness in the resulting projective flat family
\cite[Tag~\texttt{0578}]{Stacks}.  Thus
it gives exactly the geometric reductions there, with constant Hilbert
polynomial.  The complement upstairs has image in a proper closed subset
of the original base because the normalization is finite.  Noetherian
induction on that closed subset proves the assertion.  Applying this to
the divisors in $\mathcal S$ gives the required finite union of Hilbert
schemes on $X$.

\smallskip\noindent
\emph{The incidence fields.}
In that union take a positive-dimensional irreducible component of the
closure of the points corresponding to $\mathcal S$.  After replacing it
by an integral locally closed open subset $T$, the tested points are dense
in $T$, and its universal family
\[
                            Z\subset X\times T\longrightarrow T
\]
is flat with geometrically integral fibers of dimension $d-1$.  The latter
condition is available by constructibility of geometric integrality
\cite[Tags~\texttt{0579} and~\texttt{055B}]{Stacks}: the
tested fibers are integral over the algebraically closed field $l$, and
their parameter points are dense.  Choose an integral locally closed curve
$T'\subset T$, once and for all, and put $Z'=Z\times_T T'$.
Both $Z$ and $Z'$ are integral, by flatness and geometric integrality of
their generic fibers.  Their maps to $X$ are dominant.  Indeed a proper
closed subset of $X$ can contain only finitely many distinct prime
divisors, whereas the fibers over the infinitely many distinct points of
$T$, or of $T'$, are distinct divisors.

The two projections now give
\[
\begin{array}{ccccc}
 && Z' &&\\[-2pt]
 &\swarrow&&\searrow&\\[-2pt]
 X&&&&T'
\end{array}
\qquad\qquad
\begin{array}{ccccc}
 && M=l(Z') &&\\[-2pt]
 &\nearrow&&\nwarrow&\\[-2pt]
 L=l(X)&&&&P_0=l(T').
\end{array}
\]
Since $\dim Z'=d$, the extension $M/L$ is finite.  Geometric integrality
of the generic fiber of $Z'\to T'$ says precisely that $M/P_0$ is regular.
The inclusion $Z'\subset X\times T'$ also shows that its function field
is generated by the coordinates from the two factors, proving
\eqref{eq:curve-compositum}.

\smallskip\noindent
\emph{Descent along the fixed curve.}
We next show that the chosen fields work for every class in $H$.  Fix
$g\in L^\times$ with $[g]\in H$.  Let $Y$ be the intersection of the
relative smooth locus of $Z\to T$ with the inverse image of the open
subset of $X$ on which $g$ is a unit.  Every nonempty geometric fiber of
$Y\to T$ is integral.  Its fibers over the generic points of $T$ and
$T'$ are nonempty: the corresponding incidence families dominate $X$,
and the smooth locus is dense in their geometrically integral fibers.
On $Y$ consider the finite \'etale cover defined by adjoining an
$\ell$-th root of $g$.

This cover splits on a dense set of the tested fibers.  Indeed
\eqref{eq:curve-residue-trivial} gives a rational $\ell$-th root on each
such divisor, and on its smooth unit locus the root and its inverse are
regular, by normality.  The cover therefore splits on the geometric
generic fiber over $T$ as well.  Otherwise, since $\ell$ is prime and
the constants contain $\mu_\ell$, its generic Kummer polynomial would
be irreducible over that geometric function field.  The geometric generic
cover would then be integral.  Geometric integrality spreads to a
nonempty open of $T$
\cite[Tags~\texttt{0578} and~\texttt{0559}]{Stacks}, contradicting the
dense set of split fibers.

It remains to pass this splitting to our fixed curve $T'$.  Shrinking $T$
separately for each $g$ would not justify this passage: the resulting open
could miss $T'$.  Instead, descend the geometric generic splitting to a
finite extension of $l(T)$ and let $\widetilde T\to T$ be the finite
surjective normalization in that extension.  The pullback
$\widetilde Y=Y\times_T\widetilde T$ is normal, since it is smooth over
the normal scheme $\widetilde T$ \cite[Tag~\texttt{034F}]{Stacks}.
It is also integral.  Its generic
fiber is geometrically integral, and every irreducible component dominates
$\widetilde T$: on a normal scheme the components are open, and a smooth
map has open image.  Thus there is only one component.

The pulled-back finite \'etale cover of $\widetilde Y$ is generically
split, and hence split everywhere.  To justify the last implication,
each component is finite and birational over the normal integral base;
it is therefore isomorphic to that base.  Take a geometric point of
$\widetilde T$ lying over the generic point of $T'$.  The corresponding
fiber of $\widetilde Y$ is a base change of the nonempty fiber of $Y$
over that generic point.  Its split cover shows
\begin{equation}
\label{eq:curve-geometric-root}
                            g\in(M\overline{P_0})^{\times\ell}.
\end{equation}
This reasoning applies to each $g$ separately, while $T'$, $M$, and
$P_0$ stay fixed.  Lemma~\ref{lem:constant-descent}, applied to the
regular extension $M/P_0$, now gives the asserted containment in $V_M$.
\end{proof}

\subsection{Descent to the original function field}

We apply incidence to the image of a curve subfield under $\Theta$.
The last step below uses the equality $M=LP_0$ to descend all the way to
$L$, rather than merely finding a curve field after a finite extension.

\begin{theorem}[Correspondence of curve subfields]
\label{thm:curve-matching}
For every curve subfield $E\subset K$ there is a unique curve subfield
$P\subset L$ such that $\Theta(V_E)=V_P$.  This assignment is a bijection
between the curve subfields of $K/k$ and those of $L/l$.
\end{theorem}

\begin{proof}
Fix a normal integral projective model $X\subset\Pj_l^n$ of $L/l$.
Let $E\subset K$ be a curve subfield, choose $t\in E\setminus k$, and
put $H=\Theta(V_E)$.  For $a\in k$, let $h_a=\Theta([t-a])$, and let
$\mathcal S$ be the union of their mod-$\ell$ supports on $X$.
Proposition~\ref{prop:uniform-support} bounds the degrees of these
supports, and hence the degree of every prime divisor in $\mathcal S$.

The set $\mathcal S$ is infinite.  To see this, first observe that the
kernel of all divisor-order maps on $X$ is finite dimensional.  If it
contained more than $2G_X$ independent classes, apply
Proposition~\ref{prop:finite-curve-test} to representatives of a finite
independent subset of that size.  Their independent residues on the
testing curve would all have zero point orders, contradicting
Lemma~\ref{lem:curve-orders} and its genus bound.  On the other hand,
the $[t-a]$ span an infinite-dimensional subspace of $V_E$, and therefore
the $h_a$ span an infinite-dimensional subspace of $V_L$.  If
$\mathcal S$ were finite, their order vectors would lie in a finite
dimensional space; the finite-dimensional kernel just proved would give
a contradiction.

We verify \eqref{eq:curve-residue-trivial}.  For each $D\in\mathcal S$,
some $h_a$ has nonzero order at $D$.  Match the valuation of $D$ to the
quasi-divisorial valuation $v$ of $K$ using
Proposition~\ref{prop:quasi-recovery}.  Its inertia $I_v$ has nonzero
restriction to $E$.  Consequently $vE/vk\ne0$; this group is a subgroup
of the cyclic group $vK/vk$, and the valuation transcendence-degree
inequality gives $Ev=kv$.  Restriction of $D_v$ to $E$ therefore factors
through $vE/vk$ and has dimension at most one.  Explicitly, an element
whose value is in $vk$ can be rescaled by a constant to a unit; its
residue is in $kv$, so a further constant rescaling makes it a principal
unit.  Since inertia already has nonzero restriction, the images of
$I_v$ and $D_v$ on $E$ are the same line.

Now let $g\in L^\times$ with $[g]\in H$.  Whenever $\ord_D(g)=0$,
inertia at $D$ annihilates this class.  The preceding equality of
restriction images, transported through $\Phi$, implies that decomposition
at $D$ annihilates it too.  By Lemma~\ref{lem:residue-duality}, its
residue is an $\ell$-th power in $l(D)$.  Since a fixed rational function
has nonzero integral order at only finitely many prime divisors, this
proves \eqref{eq:curve-residue-trivial}.  Lemma~\ref{lem:curve-incidence}
now supplies $M/L$ and $P_0\subset M$ with $M=LP_0$, $M/P_0$ regular,
and the image of $H$ contained in the image of $V_{P_0}$.

Choose a finite normal extension $N/L$ containing $M$, allowing
inseparability, and let $P_N$ be the relative algebraic closure of $P_0$
in $N$.  It is a curve subfield of $N/l$.  The image $H_N$ of $H$ in
$V_N$ is infinite dimensional, by Lemma~\ref{lem:kummer-kernel}, and
lies in $V_{P_N}$.  Every $L$-automorphism $\gamma$ of $N$ fixes $H_N$
pointwise.  Thus $H_N\subset V_{P_N}\cap V_{\gamma(P_N)}$, and
Lemma~\ref{lem:curve-intersection} forces $\gamma(P_N)=P_N$.

Put $P_L=P_N\cap L$.  This is relatively algebraically closed in $L$.
It also has transcendence degree one over $l$.  Indeed the invariant
field $P_N^{\Aut(N/L)}$ has transcendence degree one, and it is contained
in the fixed field $N^{\Aut(N/L)}$.  For a finite normal extension the
latter is $L$ in characteristic zero and is a finite purely inseparable
extension of $L$ in positive characteristic.  In the latter case, a
uniform $p$-power of the invariant field lies in $P_N\cap L$, preserving
transcendence degree.  Thus $P_L$ is a curve subfield of $L$.

We now use the compositum identity to obtain containment already in
$V_L$.  Since $P_L\subset L\subset M$ and every element of
$P_L\subset P_N$ is algebraic over $P_0$, relative algebraic closedness
of $P_0$ in $M$ gives
$P_L\subset P_0$.  The extension $P_0/P_L$ is therefore finite algebraic.
Choose the algebraic closures inside a common algebraically closed
overfield.  Then
\begin{equation}
\label{eq:curve-exact-descent}
                          M\overline{P_0}=L\overline{P_L}.
\end{equation}
Here we used both $M=LP_0$ and the fact that $P_0$ and $P_L$ have the
same algebraic closure in that overfield.

The extension $L/P_L$ is regular.  Relative algebraic closedness is
already known, and separability is automatic in characteristic zero.
In characteristic $p>0$, choose a separating variable $u$ for $P_L/l$.
It cannot become a $p$-th power in $L$: a $p$-th root would be algebraic
over $P_L$ and hence belong to $P_L$, contrary to the choice of $u$.
Over the perfect field $l$, such an element extends to a separating
transcendence basis of $L/l$, by the usual $p$-basis criterion.  Thus
$L/l(u)$, and therefore $L/P_L$, is separable.

For each $[g]\in H$, Lemma~\ref{lem:curve-incidence} and
\eqref{eq:curve-exact-descent} give an $\ell$-th root of $g$ in
$L\overline{P_L}$.  Applying Lemma~\ref{lem:constant-descent} once more,
now to $L/P_L$, yields the exact inclusion $\Theta(V_E)\subset V_{P_L}$.

Apply the same argument to $\Theta^{-1}$ and $P_L$.  For a curve
subfield $E_1$ of $K$, it gives
\[
                       \Theta(V_E)\subset V_{P_L}\subset\Theta(V_{E_1}).
\]
Lemma~\ref{lem:curve-intersection} and the infinite dimension of $V_E$
force $E=E_1$.  Both inclusions are consequently equalities.  The same
argument with $K$ and $L$ reversed proves surjectivity of the matching,
and Lemma~\ref{lem:curve-intersection} proves its uniqueness.
\end{proof}
